\section{Calculation of $\lambda(\xi_0)$} \label{lambdacalcl}

\subsection{Series Expansion} 

The bounce-averaging normalization coefficient is defined in the case of
circular concentric flux-surfaces as (\ref{lambdacirc})
\begin{equation}
\lambda (\xi _{0})=\int_{-\theta _{c}}^{\theta _{c}}\frac{d\theta }{2\pi }%
\frac{\xi _{0}}{\xi }  \label{lambdacirc1}
\end{equation}
with
\begin{equation}
\theta _{c}=\left\{
\begin{array}{ll}
\pi  & \text{for circulating particles} \\
\theta _{T}=\cos ^{-1}\left( 1-2\xi _{0}^{2}/\xi _{0T}^{2}\right)  & \text{
for trapped particles}
\end{array}
\right.   \label{thetac}
\end{equation}

We use the expression of $\xi$ from (\ref{xi}):
\begin{equation}  \label{xi3}
\xi = \sigma\sqrt{1-\Psi(1-\xi_0^2)}
\end{equation}
where $\Psi$ has a simple expression for the case of circular flux-surface,
given in (\ref{psi}):
\begin{equation}  \label{psi1}
\Psi(\theta) = \frac{1+\epsilon}{1+\epsilon \cos(\theta)}
\end{equation}

Defining $\alpha = \theta/2$, we have $\cos(\theta) = 1-2\sin^2(\alpha)$ and
then (\ref{psi1}) becomes
\begin{eqnarray}  \label{psi2}
\Psi(\theta) & = & \frac{1+\epsilon}{1+\epsilon -2\epsilon\sin^2(\alpha)}
\nonumber \\
& = & \frac{1}{1-\xi_{0T}^2\sin^2(\alpha)}
\end{eqnarray}
where $\xi_{0T}$ is the pitch-angle at the trapped passing boundary, defined
in (\ref{xi0t}).

Then, with $\alpha_c = \theta_c/2$, (\ref{lambdacirc1}) becomes
\begin{eqnarray}  \label{lambda2}
\lambda(\xi_0) & = & \int_{-\theta_c}^{\theta_c} \frac{d\theta}{2\pi} \frac{%
\xi_0}{\sigma\sqrt{1-\Psi(\theta)(1-\xi_0^2)}}  \nonumber \\
& = & \frac{1}{\pi} \int_{0}^{\theta_c} d\theta \frac{|\xi_0|}{\sqrt{
1-(1-\xi_0^2)/(1-\xi_{0T}^2\sin^2(\alpha))}}  \nonumber \\
& = & \frac{2}{\pi} \int_{0}^{\alpha_c} d\alpha \frac{\sqrt{1-\xi_{0T}^2
\sin^2(\alpha)}}{\sqrt{1-\xi_{0T}^2/\xi_0^2\ \sin^2(\alpha)}}
\end{eqnarray}

We have the following series expansion, valid for $|x|<1$:
\begin{equation}  \label{sqrtseries}
\sqrt{1-x} = \sum_{m=0}^{\infty} \chi_m x^m
\end{equation}
with
\begin{equation}
\chi_0 = 1,\qquad \chi_m = \frac{2m-3}{2m} \chi_{m-1}
\end{equation}

Applying (\ref{sqrtseries}) to $\sqrt{1-\xi _{0T}^{2}\sin ^{2}(\alpha )}$, (%
\ref{lambda2}) becomes:
\begin{equation}
\lambda (\xi _{0})=\frac{2}{\pi }\sum_{m=0}^{\infty }\chi _{m}\xi
_{0T}^{2m}J_{2m}  \label{lambdacalc}
\end{equation}
with
\begin{equation}
J_{2m}=\int_{0}^{\alpha _{c}}d\alpha \frac{\sin ^{2m}(\alpha )}{\sqrt{1-\xi
_{0T}^{2}/\xi _{0}^{2}\ \sin ^{2}(\alpha )}}  \label{j2m}
\end{equation}

\subsection{Calculation of the Integrals $J_{2m}$}

The integrals $J_{2m}$ can be expressed in terms of complete ellipstic
integrals of the first and second kind \cite{mco86}:
\begin{eqnarray}  \label{elliptic}
K(k^2) & = & \int_0^{\pi/2} \frac{d\alpha}{\sqrt{1-k^2\sin^2(\alpha)}} =
\int_0^1 \frac{dx}{\sqrt{(1-x^2)(1-k^2 x^2)}}  \nonumber \\
E(k^2) & = & \int_0^{\pi/2} d\alpha\ \sqrt{1-k^2\sin^2(\alpha)} = \int_0^1
dx\ \frac{\sqrt{(1-k^2 x^2)}}{\sqrt{(1-x^2)}}
\end{eqnarray}

\begin{itemize}
\item  \textbf{For Trapped particles ($|\xi _{0}|<\xi _{0T}$)}, we have from
(\ref{thetac}):
\begin{equation}
\sin (\alpha _{c})=\sin (\theta _{T}/2)=\sqrt{\frac{1-\cos (\theta _{T})}{2}}%
=\frac{|\xi _{0}|}{\xi _{0T}}
\end{equation}
so that
\begin{eqnarray}
J_{0} &=&\int_{0}^{\alpha _{c}}\frac{d\alpha }{\sqrt{1-\xi _{0T}^{2}/\xi
_{0}^{2}\ \sin ^{2}(\alpha )}}  \nonumber \\
&=&\int_{0}^{\sin (\alpha _{c})}\frac{dx}{\sqrt{1-x^{2}}\sqrt{1-\xi
_{0T}^{2}/\xi _{0}^{2}\ x^{2}}}  \nonumber \\
&=&\frac{|\xi _{0}|}{\xi _{0T}}\int_{0}^{1}\frac{dx}{\sqrt{1-\xi
_{0}^{2}/\xi _{0T}^{2}\ x^{2}}\sqrt{1-x^{2}}}  \nonumber \\
&=&\frac{|\xi _{0}|}{\xi _{0T}}K\left( \frac{\xi _{0}^{2}}{\xi _{0T}^{2}}%
\right)
\end{eqnarray}
and
\begin{eqnarray}
J_{2} &=&\int_{0}^{\alpha _{c}}d\alpha \ \frac{\sin ^{2}(\alpha )}{\sqrt{%
1-\xi _{0T}^{2}/\xi _{0}^{2}\ \sin ^{2}(\alpha )}}  \nonumber \\
&=&\int_{0}^{\sin (\alpha _{c})}dx\ \frac{x^{2}}{\sqrt{1-x^{2}}\sqrt{1-\xi
_{0T}^{2}/\xi _{0}^{2}\ x^{2}}}  \nonumber \\
&=&\frac{|\xi _{0}|}{\xi _{0T}}\int_{0}^{1}dx\ \frac{\xi _{0}^{2}/\xi
_{0T}^{2}\ x^{2}}{\sqrt{1-\xi _{0}^{2}/\xi _{0T}^{2}\ x^{2}}\sqrt{1-x^{2}}}
\nonumber \\
&=&\frac{|\xi _{0}|}{\xi _{0T}}\left( \int_{0}^{1}dx\ \frac{1}{\sqrt{1-\xi
_{0}^{2}/\xi _{0T}^{2}\ x^{2}}\sqrt{1-x^{2}}}-\int_{0}^{1}dx\frac{1-\xi
_{0}^{2}/\xi _{0T}^{2}\ x^{2}}{\sqrt{1-\xi _{0}^{2}/\xi _{0T}^{2}\ x^{2}}%
\sqrt{1-x^{2}}}\right)   \nonumber \\
&=&\frac{|\xi _{0}|}{\xi _{0T}}\left[ K\left( \frac{\xi _{0}^{2}}{\xi
_{0T}^{2}}\right) -E\left( \frac{\xi _{0}^{2}}{\xi _{0T}^{2}}\right) \right]
\end{eqnarray}

\item  \textbf{For Circulating particles ($|\xi _{0}|>\xi _{0T}$)}, we have
\begin{equation}
\alpha _{c}=\frac{\theta _{c}}{2}=\frac{\pi }{2}
\end{equation}
so that
\begin{eqnarray}
J_{0} &=&\int_{0}^{\pi /2}\frac{d\alpha }{\sqrt{1-\xi _{0T}^{2}/\xi
_{0}^{2}\ \sin ^{2}(\alpha )}}  \nonumber \\
&=&\int_{0}^{1}\frac{dx}{\sqrt{1-x^{2}}\sqrt{1-\xi _{0T}^{2}/\xi _{0}^{2}\
x^{2}}}  \nonumber \\
&=&K\left( \frac{\xi _{0T}^{2}}{\xi _{0}^{2}}\right)
\end{eqnarray}
and
\begin{eqnarray}
J_{2} &=&\int_{0}^{\pi /2}d\alpha \ \frac{\sin ^{2}(\alpha )}{\sqrt{1-\xi
_{0T}^{2}/\xi _{0}^{2}\ \sin ^{2}(\alpha )}}  \nonumber \\
&=&\int_{0}^{1}dx\ \frac{x^{2}}{\sqrt{1-x^{2}}\sqrt{1-\xi _{0T}^{2}/\xi
_{0}^{2}\ x^{2}}}  \nonumber \\
&=&\frac{\xi _{0}^{2}}{\xi _{0T}^{2}}\int_{0}^{1}dx\ \frac{\xi _{0T}^{2}/\xi
_{0}^{2}\ x^{2}}{\sqrt{1-x^{2}}\sqrt{1-\xi _{0T}^{2}/\xi _{0}^{2}\ x^{2}}}
\nonumber \\
&=&\frac{\xi _{0}^{2}}{\xi _{0T}^{2}}\left( \int_{0}^{1}dx\ \frac{1}{\sqrt{%
1-x^{2}}\sqrt{1-\xi _{0T}^{2}/\xi _{0}^{2}\ x^{2}}}-\int_{0}^{1}dx\frac{%
1-\xi _{0T}^{2}/\xi _{0}^{2}\ x^{2}}{\sqrt{1-x^{2}}\sqrt{1-\xi _{0T}^{2}/\xi
_{0}^{2}\ x^{2}}}\right)   \nonumber \\
&=&\frac{\xi _{0}^{2}}{\xi _{0T}^{2}}\left[ K\left( \frac{\xi _{0T}^{2}}{\xi
_{0}^{2}}\right) -E\left( \frac{\xi _{0T}^{2}}{\xi _{0}^{2}}\right) \right]
\end{eqnarray}

\item  \textbf{For all particles}, the integrals $J_{2m},m>1$ can be derived
from $J_{0}$ and $J_{2}$ from the following recursive relation \cite{mco86}:
\begin{equation}
J_{2m}=\frac{(2m-2)}{(2m-1)}\left( 1+\frac{\xi _{0}^{2}}{\xi _{0T}^{2}}%
\right) J_{2m-2}-\frac{(2m-3)}{(2m-1)}\frac{\xi _{0}^{2}}{\xi _{0T}^{2}}%
J_{2m-4}
\end{equation}
\end{itemize}

\subsection{Truncated Expression}

It can be shown (See Figure \ref{lambdaplot}) that keeping only the two
first terms in the expansion is sufficient. Therefore $\lambda(\xi_0)$ can
be evaluated with an excellent accuracy using:
\begin{equation}  \label{lambda4}
\lambda(\xi_0) = \frac{2}{\pi} \left(J_0 - \frac{1}{2} \xi_{0T}^{2}
J_2\right)
\end{equation}
with $J_0$ and $J_2$ calculated above.

\section{Calculation of $s^{\ast}$}

The coefficient $s^{*}$ is defined for circular concentric flux-surfaces as
\begin{equation}
s^{*}=\int_{-\pi }^{\pi }\frac{d\theta }{2\pi }\frac{R\left( r,0\right) }{%
R\left( r,\theta \right) }=\int_{-\pi }^{\pi }\frac{d\theta }{2\pi }\frac{%
1+\epsilon }{1+\epsilon \cos (\theta )}  \label{sstar1}
\end{equation}

Defining
\begin{equation}
z = e^{i\theta}
\end{equation}
we get
\begin{equation}
d\theta = \frac{1}{iz} dz, \qquad cos(\theta) = \frac{1}{2}(z+1/z)
\end{equation}

and (\ref{sstar1}) becomes
\begin{equation}
s^{*}=\frac{2}{\epsilon }\frac{(1+\epsilon )}{2\pi i}\int_{\gamma }\frac{dz}{%
z^{2}+2z/\epsilon +1}  \label{sstar2}
\end{equation}
where the contour $\gamma $ is the unity circle in the complex plane.

The two roots of the denominator in (\ref{sstar2}) are
\begin{equation}
z_{\pm} = -\frac{1}{\epsilon}(1\mp\sqrt{1-\epsilon^2})
\end{equation}

Only the greater root $z_{+}$ lies inside the contour $\gamma $, so that (%
\ref{sstar2}) becomes
\begin{equation}
s^{*}=\frac{2}{\epsilon }\frac{(1+\epsilon )}{2\pi i}(2\pi i)\text{Res}%
\left( \frac{1}{z^{2}+2z/\epsilon +1},z=z_{+}\right)   \label{sstar3}
\end{equation}
where
\begin{equation}
\text{Res}\left( \frac{1}{z^{2}+2z/\epsilon +1},z=z_{+}\right) =\frac{1}{%
z_{+}-z_{-}}=\frac{1}{2/\epsilon \ \sqrt{1-\epsilon ^{2}}}  \label{sstar4}
\end{equation}
is the residue of the integral in (\ref{sstar2}) at $z=z_{+}$.

Finally, inserting (\ref{sstar4}) in (\ref{sstar3}), we find
\begin{equation}
s^{*}=H(|\xi _{0}|>\xi _{0T})\ \sqrt{\frac{1+\epsilon }{1-\epsilon }}
\label{sstar}
\end{equation}
